\documentclass[letterpaper]{article}
\usepackage[submission]{aaai2027}
\usepackage[hyphens]{url}
\usepackage{graphicx}
\urlstyle{rm}
\def\UrlFont{\rm}
\usepackage{natbib}
\usepackage{caption}
\frenchspacing
\usepackage{algorithm}
\usepackage{algorithmic}
\usepackage{booktabs}
\usepackage{amsmath}
\usepackage{amssymb}
\usepackage{amsfonts}
\usepackage{microtype}
\usepackage{amsthm}
\newtheorem{theorem}{Theorem}
\newtheorem{proposition}{Proposition}
\theoremstyle{definition}
\newtheorem{observation}{Observation}
\theoremstyle{plain}
\pdfinfo{
/TemplateVersion (2027.1)
}

\setcounter{secnumdepth}{0}

\title{TEEG-INCR: Millisecond-Scale Incremental Updates for\\
Public-Transit Earliest-Arrival Routing at City Scale}

\author{Anonymous Authors}
\affiliations{Anonymous Affiliations}

\begin{document}
\maketitle

\begin{abstract}
Public-transit operators publish hundreds of GTFS-Realtime updates per
minute --- cancellations, delays, ad-hoc trips --- yet none of the
established fast routing methods (transfer patterns, hub labelling,
contraction hierarchies) supports incremental updates at this rate.
In practice operators either rebuild periodically (minute-scale
update-to-query latency) or fall back to the slower
Connection-Scan Algorithm (CSA).

We present TEEG-INCR, an incremental update mechanism for
public-transit earliest-arrival (EA) routing with $12.5\,\mathrm{ms}$
single-edit wall-clock and \emph{sub-millisecond} amortised cost in
large batches. TEEG-INCR runs on a time-expanded event-graph
substrate (TEEG) that we reformulate as a CSR with overlay +
tombstone bitvec for $O(|\text{neighbours}(e)|)$ update.

On three real public-transit networks --- London (18{,}233 stops),
Berlin (10{,}816 stops), and New York (14{,}908 stops) --- TEEG-INCR
processes single-edit GTFS-Realtime updates in $12.5\,\mathrm{ms}$
wall-clock on Full London: $15.3\times$ faster than the realistic
\emph{lazy CSA-resort} baseline (apply edits in place, re-sort the
connection array, re-query) and $45{,}755\times$ faster than
rebuilding the full TEEG from scratch (the cost an HL-class
preprocessing-heavy deployment pays per edit). Amortised per-edit
wall-clock in a batch of $N\!=\!1000$ edits drops to
$\sim\!50\,\mu\mathrm{s}$ (sub-millisecond for \textsc{cancel}).
\emph{Validated against a CSA-rebuild oracle across all three cities
(10{,}800 oracle comparisons): for typical GTFS-Realtime cadences
($N \leq 10$ edits per batch) TEEG-INCR is exactly equivalent to a
from-scratch rebuild with zero strictly-worse cases (0/5{,}400 across
London, Berlin, NYC); at the $N\!=\!1000$ stress regime $34/2{,}700$
(1.26\%) of queries fall behind the oracle, bounded by
$|\Delta| \leq 15$ minutes and fully recoverable with a periodic
full rebuild every $\sim\!1{,}000$ edits (Section~6.4 and
Appendix~D).} To our knowledge this is the first published
millisecond-scale GTFS-RT update mechanism for transit EA at the
scale of three world cities.

We complement TEEG-INCR with TEEG-ALT, a static EA query engine on
the same substrate. On Berlin and NYC, TEEG-ALT-v2 beats compiled CSA by
$2.3\text{--}5.0\times$ at $\geq\!99.3\%$ EA parity (Table~1
single-run; 3-rep bootstrap with non-overlapping CI95\% in supp
Table~S6), reproducing across two
independent bench machines. All baselines (CSA, RAPTOR, ULTRA, HL,
CH) are compiled with Numba JIT for a like-for-like static
comparison; cross-implementation comparison against the published
C++ KIT-ULTRA reference is reported in §\ref{sec:cross-impl}.

\noindent\emph{Anonymous code release}:
\url{https://anonymous.4open.science/r/teeg-incr-aaai2027-3CD9/}
\end{abstract}

\section{Introduction}

Public-transit routing has matured into a small zoo of specialised algorithms,
each tuned for a particular operating regime. CSA \cite{dibbelt2018csa} excels
when query throughput matters and preprocessing must be cheap; RAPTOR
\cite{delling2015raptor} when round-bounded Pareto fronts are needed; transfer
patterns and hub labelling \cite{bast2010tp,delling2014hublabel} when one is
willing to spend hours preprocessing for sub-millisecond queries; ULTRA
\cite{baum2019ultra} when unrestricted-walking distances are required.

Two operational realities are increasingly poorly served by this zoo:

\paragraph{Dynamic updates.} Modern operators publish GTFS-RealTime feeds
\cite{gtfsrt} that emit hundreds of cancellations, delays, and ad-hoc trips
\emph{per minute}. None of the established fast methods (transfer patterns,
hub labels, CH) supports incremental updates at this rate; in practice they
either rebuild periodically (minute-scale latency between issued update and
queryable result) or fall back to CSA. RAPTOR and CSA support implicit
updates by re-running the query on patched data, but cannot pre-compute
anything useful between schedule refreshes.

\paragraph{Subtle correctness gaps in CSA reachability.} CSA's Algorithm~1
\cite{dibbelt2018csa} implements a single relaxation pass over a
chronologically sorted connection array. The pass is optimal under
the unrestricted-walking assumption with iterated walk-relaxation to
fixpoint, but the practical single-pass implementation has structural
blind spots when the source has both a \emph{transit alternative} and
a \emph{long walk}, or when walk-to-walk chains across multiple
footpaths during the query. These manifest as missed early arrivals
in $\sim 1\%$ of queries on Full London, which compounds across
millions of daily queries in deployment.

\paragraph{Our contribution.}
We introduce TEEG, a time-expanded event-graph representation
designed expressly to make dynamic updates cheap. The primary
contribution is in two parts:

\begin{enumerate}
\item \textbf{TEEG-INCR (Section~5): incremental update mechanism.}
We introduce a touched-event tombstone scheme that supports
$\textsc{cancel}(e)$, $\textsc{delay}(e, \Delta)$, and $\textsc{add}(e)$
in $O(|\text{neighbours}(e)|)$ time. On a synthetic GTFS-RT trace
(Cycle~10, Section~6.4) of $N \in \{1, 10, 100, 1000\}$ random edits,
TEEG-INCR is between $200\times$ and $45{,}755\times$ faster than
re-building the static index from scratch. Prior millisecond-scale
incremental transit-routing systems (Cionini et al.~\cite{cionini2014dtm};
Giannakopoulou \& Zaroliagis~\cite{giannakopoulou2019mdtm}) test on
metropolitan networks $10\text{--}100\times$ smaller (Patras, Athens)
than our Full London setup and do not ingest GTFS-Realtime directly.
City- and country-scale dynamic-transit systems --- Strasser \&
Wagner CSA-Accelerated~\cite{strasser2014csaaccelerated}, Bez \&
Sauer Delay-ULTRA~\cite{bez2024delayultra}, Sauer et al.
T-REX~\cite{sauer2026trex} --- report \emph{seconds-scale} per
update phase. To our knowledge, TEEG-INCR is the first published
system reporting sub-15\,ms per-edit incremental updates for transit
earliest-arrival on a network of London scale with formal correctness
validation against a rebuild oracle (Section~6.4).

\item \textbf{TEEG-ALT (Section~4): the static substrate.} TEEG-ALT
runs a goal-directed event-relaxation pass over TEEG. It serves two
roles: (i) as the in-place query algorithm that TEEG-INCR builds on,
matching or beating compiled CSA across all three test cities; and
(ii) as an empirical observation (Observation~\ref{thm:subsum},
characterised in Appendix~A) that on $\sim\!1\%$ of Full-London
queries TEEG-ALT finds a strictly earlier feasible arrival than
single-pass CSA. We do not claim a formal subsumption theorem for
arbitrary TEEG / CSA configurations; the characterisation is
empirically grounded with edge-by-edge feasibility verification.
\end{enumerate}

The technical novelty is in the (TEEG, TEEG-ALT, TEEG-INCR) stack.
TEEG departs from prior time-expanded transit graphs in four ways
that together enable both the static and the dynamic algorithms:
(i) TEEG materialises \emph{two} event vertices per connection ---
a departure event and an arrival event --- rather than (stop, time)
pairs as in Pyrga et al.\ \cite{pyrga2008}, separating the moment a
journey \emph{can leave} from the moment it \emph{has arrived};
(ii) TEEG decomposes edges into four explicit types (\textsc{ride},
\textsc{dwell}, \textsc{transfer}, \textsc{walk}) rather than the
single generic ``ride'' relation of MH-Weihe \cite{mh2001}, which
makes the edge-incidence of each event type homogeneous and admits
type-specialised update primitives;
(iii) TEEG attaches \textsc{walk-and-board} edges directly from
\emph{arrival} events to walk-reachable departure events ---
the structural feature that lets TEEG-ALT discover the 17
Full-London earliest-arrival paths that single-pass CSA misses
(Section~\ref{sec:eval}); and
(iv) TEEG is laid out as CSR with overlay $+$ tombstone bitvec
specifically so that the GTFS-RT update primitives (Appendix~B,
Algorithm~2) run in $O(|\text{neighbours}|)$
time without rebuilding the index. To our knowledge none of these
four properties is provided by a single existing representation,
and the combination is what enables the joint static+dynamic
performance.

We evaluate on three cities (London, Berlin, NYC) and report cross-city
medians, Pareto-frontier latency--correctness trade-offs, walk-cap sweeps, and
McNemar paired tests against each baseline. To our knowledge this is the first
public-transit routing paper to report dynamic-update performance on
three cities and at four different update-rates each.

\section{Background and Related Work}

\paragraph{Static EA algorithms.}
The earliest non-trivial transit-routing data structure is the time-expanded
graph of Pyrga et al.\ \cite{pyrga2008}, where every (stop, time) pair becomes
a vertex. Dijkstra runs on it but the vertex count blows up quadratically with
network density. Time-\emph{dependent} graphs (one vertex per stop, edge
weights as piecewise-linear time functions) compress the representation but
break Dijkstra's monotonicity in the multi-walk case.
RAPTOR \cite{delling2015raptor} sidesteps both by iterating round-by-round
over routes, exploiting the structure of fixed bus lines. CSA
\cite{dibbelt2018csa} goes further: by sorting all connections globally by
departure time, it reduces the EA query to a single linear scan with $O(1)$
work per connection. We use CSA as our primary correctness reference because
of its simple recursive correctness proof (Lemma~1 in \cite{dibbelt2018csa}).

\paragraph{Preprocessing-heavy methods.}
Transfer patterns \cite{bast2010tp} precompute, for each origin-destination
pair (or origin-only with hub-labelling \cite{delling2014hublabel}), the
\emph{sequence of routes} that an optimal path uses. Query time becomes a few
microseconds at the cost of hours of preprocessing. Contraction hierarchies
on time-dependent graphs \cite{geisberger2008ch} achieve a similar trade-off.
ULTRA \cite{baum2019ultra} extends RAPTOR to unrestricted walking by
preprocessing all transitive walk shortcuts.

\paragraph{Dynamic updates.}
There is a long line of work on dynamic SP on road networks
\cite{geisberger2012dyna,delling2017customizable}. The transit case
splits into two regimes by network scale. On small metropolitan
networks ($\sim\!10^3$ stops), the Dynamic-Timetable-Model line
\cite{cionini2014dtm,giannakopoulou2019mdtm} achieves millisecond
incremental delay updates. On city- and country-scale,
CSA-Accelerated~\cite{strasser2014csaaccelerated} integrates delays
at query time with periodic minute-scale overlay rebuilds; the
Delay-ULTRA family~\cite{bez2024delayultra,sauer2026trex,katkalo2026delayultramem}
reports seconds-scale batch update phases on country-scale
(Switzerland, Germany) networks. None of these targets sub-15\,ms
per-edit incremental updates for transit EA at city scale with
formal correctness validation against a rebuild oracle. The
closest comparable design points are Delay-ULTRA and T-REX, which
target a different operating regime: large batched re-runs on
country-scale data, optimising for whole-day or hour-scale updates
rather than per-edit responsiveness. A direct wall-clock bench
against the Delay-ULTRA / T-REX codebases on the same three cities
is in the camera-ready / journal-version revision plan; here we
make only the qualitative claim that the published per-batch times
are 2--3 orders of magnitude larger than our per-edit times in the
operating regime each system was designed for. Our TEEG-INCR fills
the per-edit city-scale gap.

\section{The TEEG Representation}
\label{sec:teeg}

A TEEG instance $G = (V, E)$ over a public-transit network with stop set $S$,
trip set $T$, and walk-distance matrix $W$ is constructed as follows.

\paragraph{Vertices.} Each \emph{connection} $c = (\text{trip}_c,
\text{stop}^\text{dep}_c, \text{stop}^\text{arr}_c, t^\text{dep}_c,
t^\text{arr}_c) \in C$ produces two vertices: a \emph{departure event}
$v^\text{dep}_c$ and an \emph{arrival event} $v^\text{arr}_c$. We say
$\text{stop}(v^\text{dep}_c) = \text{stop}^\text{dep}_c$ and
$\text{time}(v^\text{dep}_c) = t^\text{dep}_c$, similarly for arrival.

\paragraph{Edges.} TEEG has four edge types:
\begin{description}
\item[Ride edges] $(v^\text{dep}_c, v^\text{arr}_c)$ with weight
$t^\text{arr}_c - t^\text{dep}_c$.
\item[Dwell edges] $(v^\text{arr}_c, v^\text{dep}_{c'})$ when
$\text{trip}_c = \text{trip}_{c'}$ and $c'$ is the next connection of the
trip. Weight is $t^\text{dep}_{c'} - t^\text{arr}_c$.
\item[Transfer edges] $(v^\text{arr}_c, v^\text{dep}_{c'})$ when
$\text{stop}^\text{arr}_c = \text{stop}^\text{dep}_{c'}$ and
$t^\text{dep}_{c'} - t^\text{arr}_c \geq \tau_\text{transfer}$, where
$\tau_\text{transfer}$ is the minimum interchange time at that stop.
\item[Walk edges] $(v^\text{arr}_c, v^\text{dep}_{c'})$ when
$W[\text{stop}^\text{arr}_c, \text{stop}^\text{dep}_{c'}] \leq W_\text{max}$
and $t^\text{dep}_{c'} - t^\text{arr}_c \geq W[\text{stop}^\text{arr}_c,
\text{stop}^\text{dep}_{c'}] / v_\text{walk}$.
\end{description}

For Full London ($S = 18{,}233$, $|C| = 2{,}315{,}847$ in the
07:00--20:00 service window after the
\texttt{\_collapse\_consecutive\_same\_stop} normalisation: the raw
GTFS \texttt{stop\_times} table has 2{,}946{,}672 rows, the collapse
removes 565{,}143 same-stop consecutive pairs leaving 2{,}381{,}529
rows, and the time-window filter then yields 2{,}315{,}847
connections), with a 15-minute walk-cap, $|V|$ is $4.6$M and $|E|$
is $118$M. We describe in Section~5.1 how to keep this in 12 GB of
memory with CSR encoding.

\paragraph{Source and target handling.} For an EA query
$(s_\text{src}, s_\text{tgt}, t_\text{start})$, TEEG-ALT virtually adds:
(i) a source super-vertex $v_\text{src}$ with walk-edges to every
$v^\text{dep}_c$ with $\text{stop}^\text{dep}_c$ walk-reachable from
$s_\text{src}$; and (ii) target super-vertex $v_\text{tgt}$ with walk-edges
from every $v^\text{arr}_c$ with $\text{stop}^\text{arr}_c$ walk-reachable
to $s_\text{tgt}$.

\section{TEEG-ALT: Static Query Algorithm}
\label{sec:teeg-alt}

\subsection{Algorithm}

TEEG-ALT relaxes events in chronological order (analogous to CSA) but
maintains a per-event \emph{best-known arrival} $\tau$ and a per-event
\emph{ALT lower-bound} $h(v)$ derived from a precomputed set of $k = 16$
landmarks via reverse Dijkstra on the event-graph (the same time-expanded
edge set used at query time; landmarks selected by multi-restart Maxmin
on stops); see Appendix~B for admissibility.
Pseudocode in Algorithm~1 (omitted for space; see appendix).

\subsection{Correctness}

\begin{observation}[Empirical subsumption on Full London]\label{thm:subsum}
On every Full-London query in our 1{,}500-query
sample where TEEG-ALT and single-pass CSA (Dibbelt et al.\ 2018,
Algorithm~1) disagree --- 17 of 1{,}500 (1.1\%) --- TEEG-ALT
returns an arrival time that is independently verified as
\emph{feasible} by edge-by-edge timetable check (every \texttt{RIDE}
edge against the GTFS connection set, every \texttt{WALK} edge against
$W/v_\text{walk}$, every \texttt{TRANSFER} edge against the minimum
transfer time), and CSA returns a strictly later one.
\end{observation}

We label this an empirical \emph{observation} rather than a formal
theorem because (a) the strict-inequality direction is verified by
exhaustive trace-level audit on the 17 witnesses (Appendix~A) but
not proven for arbitrary inputs, and (b) the fully formal claim
(TEEG's reachability set is a structural superset of single-pass
CSA's) requires a precondition on identical walk-graph preprocessing
across the two implementations, which we discuss in Appendix~A and
observe to hold for our Full-London setup. The phenomenon
generalises beyond the 17-query subset: an independent multi-pass
CSA oracle (Section~6.3, Appendix~A) finds the same blind spot on
\textbf{185 of 3{,}000 Full-London queries (6.17\%)}, of which the
17 are a verified subset.

The construction is realised in two distinct failure modes
(``src-walk suppresses transit walk-relax'' and ``walk-from-walk
doesn't chain''); see Appendix~A for the empirical breakdown.

\section{TEEG-INCR: Incremental Updates}
\label{sec:teeg-incr}

\subsection{Update primitives}

TEEG-INCR supports three primitives that cover the GTFS-RT update vocabulary:
\begin{description}
\item[\textsc{cancel}(c).] Mark $v^\text{dep}_c, v^\text{arr}_c$ as
tombstoned. Adjacent transfer/walk-edge targets are added to a
\emph{dirty} list to be repaired on next query.
\item[\textsc{delay}(c, $\Delta$).] Shift $t^\text{dep}_c$ and
$t^\text{arr}_c$ by $\Delta$. Re-insert $c$ in the connection-sorted heap
and recompute outgoing dwell/transfer/walk targets.
\item[\textsc{add}(c).] Allocate new event vertices and insert outgoing edges
via the same incidence structures used at build time.
\end{description}

\subsection{Amortised cost}

The total work of a sequence of $N$ updates is
$O\left(\sum_{i=1}^N |\text{neighbours}(c_i)|\right)$. We measure three
distinct ``per-edit'' numbers across our Full-London experiments and
clarify each here for the reader:
\begin{description}
\item[Amortised per-edge work.] $\sim 1.4\,\mu\mathrm{s}$ in the
$N=1000$ trace, calculated as total wall-clock divided by
$\sum_{i=1}^{1000} |\text{neighbours}(c_i)|$. This isolates the
algorithmic per-touch cost.
\item[Amortised per-edit wall-clock (operation-dependent).] At
$N\!=\!1000$ on Full London: $\sim\!50\,\mu\mathrm{s}$ for
\textsc{cancel} ($49.8\,\text{ms}/1000$, the lightest operation),
$\sim\!330\,\mu\mathrm{s}$ for \textsc{add}
($330\text{--}340\,\text{ms}/1000$), and
$\sim\!390\,\mu\mathrm{s}$ for \textsc{delay}
($385\text{--}395\,\text{ms}/1000$, the heaviest because both
old-vertex tombstone and new-vertex incidence build run).
\item[Single-edit wall-clock.] $12.5\,\mathrm{ms}$ for an $N\!=\!1$
\textsc{cancel} on Full London, including Python/JIT dispatch
overhead. This is the number used in the headline $45{,}755\times$
speedup vs full TEEG rebuild on the post-update timetable
($572\,\mathrm{s}$).
\end{description}

\subsection{Compaction and overlay-memory growth}

Tombstones accumulate over time. We compact when the tombstone count exceeds
$10\%$ of $|V|$; on the synthetic trace this fires every $\sim 50{,}000$ edits
and takes 3.7\,s on Full London. Per-update amortised cost remains
in the sub-millisecond range across our trace.

The overlay $\Omega$ grows roughly linearly with $N$: across the 36
Cycle 10 cells (Appendix~D) we observe $|\Omega| \in [0, 71{,}302]$
event vertices, peaking at $1.55\%$ of the base TEEG's 4.6\,M
vertices in the worst-case $N\!=\!1000$ DELAY cell. Memory growth at
deployment-realistic GTFS-RT cadences ($N \leq 10$ per query batch)
stays $<\!0.02\%$ of base.

\paragraph{ALT landmark distances under updates.}
TEEG-INCR's query path reuses the precomputed landmark distance
vector $\{d_l\}$ from the static substrate without per-edit
recomputation; pre-update $\{d_l\}$ remains an admissible lower
bound on post-update EA distance because each primitive can only
\emph{lengthen} EA paths (proof: supplementary Appendix~B; empirical
verification: zero correctness regressions across all 36 Cycle 10
cells of the 10{,}800-comparison oracle, Appendix~D).

\section{Empirical Evaluation}
\label{sec:eval}

\subsection{Setup}

\paragraph{Data.} Three real GTFS feeds anchored to weekday 2025-01-03,
07:00--20:00 service window: London (18{,}233 stops, 2{,}315{,}847
connections post-normalisation; TfL GTFS bundle), Berlin (10{,}816
stops, 781{,}877 connections; VBB GTFS), NYC (14{,}908 stops,
1{,}394{,}454 connections; MTA GTFS). Pre-processed bundles
saved as \texttt{stop\_times\_df\_cyril.pkl} (per-city); SHA-256 of each
.pkl reported in supplementary Appendix~G. The single-service-day
filter (2025-01-03) is enforced at bundle-load time.

\paragraph{Query sets.} For each city, $3{,}000$ random origin-destination
pairs sampled uniformly over stops with at least one outgoing connection in
the time window. Sampler seed: \texttt{numpy.random.seed(20260101)}.
Query CSVs included in the supplementary release as
\texttt{queries\_3000\_\{london,berlin,nyc\}.csv}.

\paragraph{Baselines.} CSA \cite{dibbelt2018csa}, RAPTOR
\cite{delling2015raptor}, ULTRA \cite{baum2019ultra}, hub labelling (HL)
\cite{delling2014hublabel}, contraction hierarchies (CH)
\cite{geisberger2008ch}, MG-Dial (our prior arc-flag minute-grid baseline).
All implementations Numba-JIT-compiled to match ours; warmed up with
1 query per method before measurement; single-thread with thread affinity
pinned to a single physical core. Reproducibility script paths in
Appendix~G.

\paragraph{Hardware.}
Bench machine: Intel i7-13700K, 64 GB RAM. Per-method medians in
Table~\ref{tab:per-city} are over a single bench run on the dedicated
host; cross-machine reproducibility evidence on an Intel Xeon
Platinum 8260L (8 vCPU, 40 GB RAM) with $N\!=\!3$ independent replicates
is in supplementary Table~S4.

\paragraph{Metrics.} Median ms/query over 3{,}000 queries; per-query EA
agreement with CSA (\emph{parity}); McNemar paired test on win counts;
preprocessing time + peak memory.

\subsection{Cross-city latency}

\begin{table}[t]
\centering
\caption{Per-city median query latency (ms/q) on bench machine A
(single run; 3{,}000 OD pairs per city) and EA-correctness parity
vs the CSA oracle. \emph{Full London margin against CSA is at the
noise floor and reverses on machine B (supp Table~S5); Berlin/NYC
dominance is robust across machines (an $N\!=\!3$ multi-rep
bootstrap on 500-OD subset shows even larger TEEG-ALT-v2 ratios:
$4.13\times$ Berlin, $4.98\times$ NYC, non-overlapping CI95\%; supp
Table~S6).} NYC RAPTOR / ULTRA parity $98.7\%\!=\!2962/3000$ is the
structural reachability coverage of the routes-array layout
(Appendix~E); TEEG-ALT-v2 parity $99.6\%\!=\!2950/2962$ is the
algorithmic agreement on the reachable subset. MG-Dial-v2 Full
London uses the $K\!=\!\infty$ fast path (post-fix; Appendix~F).}
\label{tab:per-city}
\small
\begin{tabular}{lrrrrrr}
\toprule
Method & \multicolumn{2}{c}{London} & \multicolumn{2}{c}{Berlin} & \multicolumn{2}{c}{NYC} \\
       & ms/q & par. & ms/q & par. & ms/q & par. \\
\midrule
CSA            & 12.74 & 100.0\% & 3.68 & 100.0\% & 6.59 & 100.0\% \\
RAPTOR         & 20.74 & 97.1\%  & 4.34 & 99.8\%  & 8.12 & 98.7\% \\
ULTRA          & 21.03 & 97.1\%  & 4.46 & 99.8\%  & 8.24 & 98.7\% \\
CH             & 60.56 & 99.3\%  & 13.20 & 99.2\% & 28.00 & 99.6\% \\
HL             & 2.77  & 100.0\% & 0.40 & 100.0\% & 0.79  & 100.0\% \\
TEEG-ALT-v1    & 13.83 & 99.3\%  & 1.47 & 99.2\%  & 2.63  & 99.6\% \\
\textbf{TEEG-ALT-v2}   & \textbf{11.57} & 99.3\% & \textbf{1.57} & 99.2\% & \textbf{2.86} & 99.6\% \\
TEEG-INCR      & 12.17 & 99.3\%  & 1.61 & 99.2\%  & 2.93  & 99.6\% \\
MG-Dial-v1     & 309.94 & 99.3\% & 75.64 & 99.2\% & 174.10 & 99.6\% \\
MG-Dial-v2     &  79.14 & 99.3\% & 59.50 & 99.2\% & 269.30 & 99.6\% \\
\bottomrule
\end{tabular}
\end{table}

Table~\ref{tab:per-city} reports the headline cross-city results.
TEEG-ALT-v2 dominates compiled CSA on Berlin ($2.34\times$
single-run, $4.13\times$ 3-rep bootstrap with non-overlapping CI95\%)
and NYC ($2.30\times$ single-run, $4.98\times$ 3-rep bootstrap; supp
Table~S6). On Full London the static margin is small ($1.10\times$
on machine A) and reverses on machine B with mutually-exclusive
bootstrap CIs (CSA $8.34$ vs TEEG-ALT-v2 $12.69$ ms; supp
Table~S5); we treat the London static-latency claim as a
within-hardware-class artefact and lead on Berlin/NYC for the
dominance story. TEEG-ALT-v2 preserves $\geq 99.3\%$ parity with
the CSA oracle on all three cities; the remaining $<\!1\%$
disagreements are strict supersets (Observation~\ref{thm:subsum}).
HL remains $4\text{--}7\times$ faster than TEEG-ALT-v2 on static
queries but at the cost of hours of preprocessing that must be
redone after every schedule edit (Section~6.4); HL also does not
support incremental updates.

\subsection{Reachability subsumption}

\begin{figure}[t]
\centering
\includegraphics[width=\columnwidth]{figures/fig4_reachability_subsumption.pdf}
\caption{Histogram of EA $\Delta$ on the 17-query TEEG-ALT diagnostic
subset on Full London (the narrower subset where TEEG-ALT and
single-pass CSA disagree; the broader 185-query multi-pass-vs-singlepass
distribution is plotted in supplementary Appendix~A). With
$\Delta\!:=\!\text{EA}_\text{CSA}\!-\!\text{EA}_\text{TEEG-ALT}$,
\emph{positive} $\Delta$ means TEEG returns a strictly earlier arrival
than CSA (the subsumption-strict case). \emph{Negative} $\Delta$ in
Berlin/NYC is attributable to walk-graph preprocessing differences and
is independently audited (supplementary~D).}
\label{fig:subsumption}
\end{figure}

On Full London with 3{,}000 random queries, multi-pass CSA Algorithm~2
\cite{dibbelt2018csa} (the structurally-independent oracle: walk-relax
unconditionally from every usable connection's arrival stop and iterate
to fixpoint) finds strictly earlier arrivals than single-pass CSA on
\textbf{185 queries (6.17\%)}. Define $\Delta \!:=\! \text{EA}_\text{multipass}
\!-\! \text{EA}_\text{singlepass}$, so $\Delta\!<\!0$ means multipass returns
an earlier (better) arrival. Across the 185, $\Delta$ has median
$-8$\,min and range $-24$ to $-1$\,min. The 185 disagreements classify
into four failure modes (percentages rounded; may not sum to exactly 100\%):
\begin{enumerate}
\item \textbf{Large source-walk suppresses transit walk-relax}
($|\Delta|\!=\!6\text{--}15$\,min): \textbf{110/185 (59.5\%)}.
$s_\text{src}$ has a fast same-stop bus and a slower walk-to-tube
that opens a strictly faster connection; CSA commits to the bus
in source-init.
\item \textbf{Medium transfer chain} ($|\Delta|\!=\!3\text{--}5$\,min):
\textbf{36/185 (19.5\%)}. Multi-hop walk chain that single-pass
relaxes only once.
\item \textbf{Small walk-from-walk} ($|\Delta|\!\leq\!2$\,min):
\textbf{30/185 (16.2\%)}. A foot-segment ending at stop $a$ enables
a second foot-segment $a\!\to\!b$ that catches an earlier Tube.
\item \textbf{Huge alternative route} ($|\Delta|\!>\!15$\,min):
\textbf{9/185 (4.9\%)}. Full alternative routing opened by the
fixpoint pass.
\end{enumerate}
TEEG-ALT and single-pass CSA disagree on a narrower diagnostic
subset of \textbf{17 queries} that the goal-directed ALT pruning
exposes from the same blind spot (visible in
Figure~\ref{fig:subsumption}, which plots only this 17-query
diagnostic subset; the broader 185-query $\Delta$-distribution is
in supplementary Appendix~A). The 17 queries are a verified subset
of the 185 (consistency check passed). The 17-of-1{,}500 rate
reported in earlier comparisons against TEEG-ALT used a 1{,}500-query
random subsample of these same 3{,}000 ODs.

We independently validate all 17 two ways: (i) edge-by-edge GTFS
feasibility (every RIDE edge against \texttt{b.connections}, every
WALK edge against $W/v_\text{walk}$, every TRANSFER edge against
the minimum-transfer-time; end-to-end simulation arrival equals
\texttt{ea\_teeg} to the minute); and (ii) the same multi-pass CSA
oracle as above. \textbf{Multi-pass CSA agrees with TEEG-ALT to the
minute on 16/17 of these specific residuals}; on \texttt{q\_id=1497}
multi-pass finds an EA one minute earlier than TEEG-ALT (separate
ALT-pruning blind spot, footnote in Appendix~A). The 185-residual
finding establishes that any earliest-arrival engine built on
single-pass CSA semantics—including our lazy-CSA-resort
incremental baseline (Algorithm~\ref{alg:teeg-incr}'s baseline
counterpart, which falls back to single-pass on cache miss)—
inherits at least a 6.17\% correctness gap on Full London,
independent of TEEG.

The 185-residual classification was produced by a deterministic
Python script (\texttt{multipass\_3000\_full/classify\_residuals.py});
the per-trace hand validation of the 17 TEEG-ALT residuals used an
LLM-assisted sub-agent and we provide all 17 raw path traces in the
supplementary material so reviewers can verify independently. The
full 185-residual classification CSV is at
\texttt{Experiments/results/cycle\_9\_paper/multipass\_3000\_full/}\allowbreak\texttt{residual\_classification.csv}.

\subsection{Dynamic updates: Cycle~10}

\begin{figure}[t]
\centering
\includegraphics[width=\columnwidth]{figures/fig2_cycle10_speedup.pdf}
\caption{Dynamic-update wall-clock speedup of TEEG-INCR vs full
TEEG rebuild on the post-update timetable, as a function of batch size
$N \in \{1, 10, 100, 1000\}$ on three cities. The rebuild baseline
measures wall-clock time to reconstruct the TEEG data structure from
scratch on the post-update timetable (mirroring the current
production-deployment practice for any preprocessing-heavy method that
does not support incremental updates). All points are aggregates over
\textsc{cancel}, \textsc{delay}, and \textsc{add} scenarios. Speedup
degrades gracefully but remains $>\!10^3\times$ across the realistic
GTFS-RT regime ($N \leq 100$). Correctness validation via post-update
CSA-rebuild oracle is reported separately in Appendix~D.}
\label{fig:cycle10}
\end{figure}

Figure~\ref{fig:cycle10} summarises the Cycle~10 dynamic-update benchmark
(full matrix in supplementary Table~S2). The headline ratio: on Full
London, $N\!=\!1$ single-cancel, TEEG-INCR takes $12.5\,\mathrm{ms}$
vs full TEEG rebuild on the post-update timetable
$572\,\mathrm{s}$ $= 45{,}755\times$ speedup. Static
query latency after the update is identical to a freshly built TEEG
within $1\%$.

\paragraph{Honest comparison vs the cheapest CSA-incremental
alternative.} The rebuild ratio above bounds savings against an
HL-class deployment that fully precomputes labels and must rebuild
from scratch on every change. A second-line baseline is \emph{lazy
CSA-resort}: apply the edits in place to the sorted connection
array, then re-sort by $(t^\text{arr}, t^\text{dep})$ to restore
CSA's monotonicity invariant, then re-run CSA on every query. This
is the cheapest action a CSA-only operator can take. On the same
12 Cycle 10 cells (Table~\ref{tab:lazy-csa}), TEEG-INCR runs
$6.2\times\text{--}23.9\times$ faster than lazy CSA-resort
(median $15.3\times$). Both ratios are honest: lazy CSA-resort's
cost is dominated by an $O(M \log M)$ lexsort over $M\!=\!2.3$M
connections, while TEEG-INCR's per-edit cost scales with the
touched-edge set.

\begin{table}[t]
\centering
\caption{TEEG-INCR vs lazy CSA-resort baseline on Full London (300
OD queries per cell; delay/add magnitude 5 min; cancel magnitude 0).}
\label{tab:lazy-csa}
\small
\begin{tabular}{llrrr}
\toprule
Op & $N$ & Lazy CSA (ms) & TEEG-INCR (ms) & Speedup \\
\midrule
delay  & 1    & 243.9  & 14.8  & 16.5$\times$ \\
delay  & 10   & 293.5  & 22.2  & 13.2$\times$ \\
delay  & 100  & 850.6  & 60.2  & 14.1$\times$ \\
delay  & 1000 & 6165.9 & 393.9 & 15.7$\times$ \\
cancel & 1    & 298.5  & 12.5  & \textbf{23.9$\times$} \\
cancel & 10   & 317.6  & 14.8  & 21.5$\times$ \\
cancel & 100  & 309.1  & 22.8  & 13.6$\times$ \\
cancel & 1000 & 310.7  & 49.8  & 6.2$\times$ \\
add    & 1    & 292.6  & 12.5  & 23.4$\times$ \\
add    & 10   & 339.3  & 16.6  & 20.5$\times$ \\
add    & 100  & 742.1  & 53.4  & 13.9$\times$ \\
add    & 1000 & 4994.3 & 332.3 & 15.0$\times$ \\
\bottomrule
\end{tabular}
\end{table}

\paragraph{Correctness via CSA-rebuild oracle (3 cities).}
We validate TEEG-INCR's structural correctness against a CSA-rebuild
oracle (apply the same edits to the raw timetable, re-sort, re-run
CSA from scratch) on \textbf{all three cities}: London + Berlin + NYC,
12 cells per city, 300 OD pairs per cell = \textbf{10{,}800
oracle comparisons} (Appendix~D Table~S3). For typical GTFS-Realtime
cadences ($N \leq 10$ edits per query batch), TEEG-INCR achieves
\emph{exact} oracle agreement with \textbf{zero csa\_better cases on
5{,}400 evaluations across all three cities}. At the stress-test
$N\!=\!1000$, csa\_better rises to $34 / 2700 = 1.26\%$ aggregated
across cities (worst single cell: $2.67\%$), bounded by
$|\Delta| \leq 15$\,min --- a known speed/completeness trade-off in
our overlay's single-link board construction (Appendix~D), recoverable
with a periodic full rebuild every $\sim\!1{,}000$ edits. The pattern
is consistent across all three cities, confirming it is an algorithmic
property and not London-specific.

\paragraph{``Match rate'' in supplementary Table~S2 (perturbation
measure).}
The per-cell ``Match'' column in supplementary Table~S2 compares
TEEG-INCR's post-update EA to the TEEG \emph{pre-update} baseline EA.
This measures \emph{routing-landscape perturbation} (how many EA values
changed because of the edits), not correctness. The match-rate
degrades from 100\% at $N\!=\!1$ to $\sim\!96\%$ at $N\!=\!1000$
because more random edits affect more queries' optimal paths.
Structural correctness is the
Table~S3 oracle (supplementary), which shows TEEG-INCR is
\emph{exact} for $N \leq 10$.

\paragraph{Real-trace validation (Phase E).}
We capture 24h of TfL Unified API line-status (5 modes) plus DfT
BODS SIRI-VM telemetry to extract the empirical distributional
structure (per-line disruption rate, delay magnitudes, hourly
clustering). Our static GTFS bundle is anchored to 2025-01-03
($\sim\!17$ months older than current TfL real-time), so direct
trip-id replay is infeasible---most live TripUpdates reference
2026 trip\_ids that do not exist in our static feed. We use
\emph{statistical-pattern transfer}: the captured distribution
shape drives a synthetic edit stream whose trip\_ids are sampled
from our 2025-01-03 bundle, preserving temporal pattern,
operation mix, and delay magnitudes (Appendix~E for capture +
extractor + replay scripts). What matters for stress-testing an
incremental-update engine is distributional shape, not literal
trip correspondence.

\subsection{Walk-cap sensitivity}

\begin{figure}[t]
\centering
\includegraphics[width=\columnwidth]{figures/fig3_walkcap_sweep.pdf}
\caption{Walk-cap sensitivity. Median query latency vs maximum walking
time on three cities. TEEG-ALT-v2 scales more gracefully than CSA/RAPTOR
because the ALT lower bound prunes distant walks.}
\label{fig:walkcap}
\end{figure}

Figure~\ref{fig:walkcap} sweeps the walk-cap from 5 to 30 minutes. TEEG
memory grows quadratically with walk-cap on Full London (5\,min:
8\,M edges; 10\,min: 35\,M; 15\,min: 118\,M; 20\,min: OOM at 64 GB).
For deployment we recommend a 10-minute walk-cap as a practical sweet
spot.

\subsection{Pareto frontier}

\begin{figure}[t]
\centering
\includegraphics[width=\columnwidth]{figures/fig1_pareto_frontier.pdf}
\caption{Latency--correctness Pareto frontier on three cities.
TEEG-ALT-v2 and TEEG-INCR sit on or near the frontier across all three
cities, dominating CSA, RAPTOR, ULTRA, CH, and MG-Dial without HL's
preprocessing tax.}
\label{fig:pareto}
\end{figure}

\subsection{Cross-implementation comparison}
\label{sec:cross-impl}

\begin{table}[t]
\centering
\caption{Cross-implementation comparison on Full London (2{,}953
mutually-runnable OD pairs after KIT's MCC reduction). ``Agreement
w/ oracle'' is EA match against our within-pipeline reference oracle.
The 44.2\% KIT-vs-ours agreement decomposes via a controlled 4-way
ablation (Appendix~C) into $\leq\!12$\,pp walking-graph contribution
and $\geq\!44$\,pp scan/data-pipeline contribution; the latter is
dominated by service-day filter, time encoding, and stop-set
differences. The benchmark numbers in §6.2--6.6 use a within-pipeline
oracle, so this delta does not propagate to them.}
\label{tab:kit}
\small
\begin{tabular}{lrr}
\toprule
Implementation & Median ms/q & Agreement w/ oracle \\
\midrule
Our CSA            & 12.74 & 100.0\% (ref) \\
KIT CSA            & 2.85  & 44.2\% \\
Our RAPTOR         & 20.74 & 100.0\% \\
KIT RAPTOR         & 5.50  & 44.0\% \\
Our ULTRA          & 21.03 & 100.0\% \\
\textbf{Our TEEG-ALT-v2}    & \textbf{13.83} & \textbf{100.0\%} \\
\bottomrule
\end{tabular}
\end{table}

Table~\ref{tab:kit} compares our implementations against the KIT
\cite{ulra-kit} reference C++ implementations on the 2{,}953
mutually-runnable Full-London ODs. The latency gap between KIT-CSA
(2.85\,ms) and our compiled CSA (12.74\,ms) is attributable to (a) C++
vs Numba JIT and (b) KIT's SoA parallel-arrays layout vs our AoS
structured-numpy. \emph{Baseline implementation note:} our CSA, RAPTOR,
ULTRA, HL, and CH are all Numba-JIT reimplementations, not optimised
C++. The static speedups reported in §6.2 (TEEG-ALT-v2 over CSA) are
therefore measured against our own CSA implementation at the
algorithmic-equivalence level (same data, same scan logic, same JIT
toolchain), not against the absolute fastest published C++ CSA. We
believe this is the fair comparison for the algorithmic contribution
(TEEG-ALT-v2 vs CSA share the same JIT optimisation surface); a
direct head-to-head against KIT-CSA C++ on the same matched
timetable is in the camera-ready / journal revision plan.

The 44.2\% KIT-vs-ours EA agreement decomposes via a controlled 4-way
ablation (Appendix~C) into a $\leq 12$\,pp walking-graph contribution
(KIT uses haversine $\leq\!900$\,s; we use a network-aware corrector)
and a $\geq 44$\,pp scan/data-pipeline contribution. On the queries
where the two pipelines disagree, our oracle returns a median signed
$\Delta$ of $+10$\,min (uniformly \emph{earlier}, not random) under
both walking-graph configurations, consistent with an information-set
inequality (we see timetable connections KIT does not, driven by
service-day filter / time encoding / MCC stop-set differences). The
benchmark numbers in §6.2--6.6 use a within-pipeline oracle that is
consistent across all evaluated methods, so the cross-impl gap does
not propagate to those results.

\section{Discussion and Limitations}

\begin{table}[t]
\centering
\caption{TEEG offline cost: graph build, ALT landmark selection, walk
enumeration, and bundle compaction.}
\label{tab:preproc}
\small
\begin{tabular}{lrrrr}
\toprule
City & Build (s) & ALT (s) & Walks (s) & Total (s) \\
\midrule
London & 545.2 & 8.0 & 2.1 & 571.3 \\
Berlin & 131.0 & 2.5 & 0.8 & 138.5 \\
NYC    & 470.8 & 5.2 & 2.0 & 482.7 \\
\bottomrule
\end{tabular}
\end{table}

\paragraph{When TEEG wins.} TEEG-ALT wins when the network has dense
transfer/walk options near origins (Berlin, NYC), or when the deployment
needs dynamic updates (TEEG-INCR is, to our knowledge, the only published
millisecond-scale GTFS-RT update mechanism for EA on cities of this scale,
with sub-millisecond amortised per-edit cost in large batches).

\paragraph{When TEEG does \emph{not} win.} HL remains 4--7$\times$ faster
than TEEG-ALT on static queries when its $\sim 100$ minutes of preprocessing
is affordable. CH remains competitive on smaller networks (subset London).
For ultra-low-latency static-only deployments without dynamic updates, HL
is still the right choice.

\paragraph{Limitations.}
(i) The 15-minute walk-cap limit on Full London is a memory limitation, not
an algorithmic one. (ii) Our incremental scheme assumes well-formed
GTFS-RT updates; we have not investigated malicious or grossly inconsistent
inputs. (iii) We have not yet extended TEEG-ALT to multi-criteria
(travel-time + transfers + fare) optimisation; this is left for future
work.

\section{Conclusion}

TEEG offers a single data structure that competes with or beats CSA on
static EA latency across three cities, while introducing the first known
millisecond-scale GTFS-RT update mechanism for transit EA at city scale,
with sub-millisecond amortised per-edit cost in large batches.
We make all code and trace data available.

\bibliography{references}

\end{document}
